% ===== Block A: PDF p.2–3 | 原文 §1 Introduction + 致谢/资助 | 唯一编号公式 1.1–1.2 =====
% [原文 p.2]

作为一个物理理论，二维引力是相当平凡的理论，因为 Einstein--Hilbert 作用量
\begin{equation}
  S = \frac{1}{2\kappa} \int_{\Sigma} \dd^2 x \, R \sqrt{-h}
  \tag{1.1}
\end{equation}
是曲面 $\Sigma$ 的一个拓扑不变量。因此，Einstein 方程自动被满足。与之相反，二维\emph{量子}引力是一个丰富得多的理论，与可积系统理论和代数几何有着深刻的联系。在量子框架下，物理上实现的并不是曲面 $\Sigma$ 上一个固定的度规 $h$，而是一个动力学地涨落的度规。我们关心的量——路径积分——于是就是在所有这类度规按对称性取商得到的空间上的积分：
% [Tutor Note: (1.2) 的记号是“集合 / 群作用”的商记号：分子是曲面与其上度规的全体，分母的两行分别是微分同胚群与共形（外尔）变换群；两者都作用在分子上，故要取商。这个商空间正是（带标记点的）黎曼曲面模空间 $\cM_{g,n}$ 的物理原型。核对方法：若不对分母中的共形变换取商，得到的是 Teichmüller 空间；再对微分同胚（模去同痕）取商才得到模空间。]
\begin{equation}
  \left\{ (\Sigma , h) \;\middle|\;
  \begin{array}{c}
    \text{曲面 } \Sigma \\[2pt]
    \text{带度规 } h
  \end{array}
  \right\}
  \Bigg/
  \begin{array}{c}
    \text{微分同胚} \\[2pt]
    \text{共形变换}
  \end{array}
  \tag{1.2}
\end{equation}
用数学的语言说，我们感兴趣的是参数化黎曼曲面的空间，更确切地说，是计算这类模空间上的积分。

% [Tutor Note: 由 Gauss--Bonnet 定理，二维中 $\int_\Sigma R\sqrt{-h}\,\dd^2x = 4\pi\chi(\Sigma) = 4\pi(2-2g)$，故 (1.1) 只依赖于亏格 $g$，与度规 $h$ 的具体形状无关——这正是原文所说“是拓扑不变量”的依据（等价于取 $\kappa=1$ 时 $S = 4\pi(2-2g)/(2\kappa)$）。也正因作用量不含度规的动力学自由度，其变分恒为零，Einstein 方程“自动满足”：在二维，Einstein 张量 $R_{\mu\nu} - \tfrac12 h_{\mu\nu} R$ 恒等于零。]

研究二维量子引力还有一条完全不同的途径，它建立在把曲面离散化并对三角剖分计数的想法之上，而这又与随机矩阵理论相关。“随机矩阵方法”始于 G.~'t~Hooft 于 1974 年 [Hoo74] 的发现：他在研究强核相互作用时发现，矩阵积分与画在曲面上的图自然地联系在一起，并按图的拓扑加权。't~Hooft 这个最初的例子，随后被物理学家 E.~Brezin、C.~Itzykson、G.~Parisi 和 J.-B.~Zuber [BIPZ78] 推广为计数映射的一般范式。用他们的方法，他们重新得到了数学家 W.~T.~Tutte 在 60 年代关于球面三角剖分个数的一些结果 [Tut68]。
% [Tutor Note: 't~Hooft 的发现可以精确表述为：矩阵积分的 $1/N$ 展开就是按亏格的展开。一个 $N\times N$ 矩阵模型的配分函数按 $N^{2-2g}$ 分块，其中 $g$ 正是所画图的亏格。这条“大 $N$ 展开 $=$ 拓扑展开”的观察，是本讲义把矩阵模型与黎曼曲面模空间联系起来的出发点。核对方法：取单矩阵高斯模型，其 $1/N$ 展开首项对应球面（平面图），次项对应环面。]

在连续极限下，人们预期这两条途径会重合。这两个二维量子引力模型等价的想法有着惊人的推论，并促使 E.~Witten 提出他关于黎曼曲面模空间几何的著名猜想 [Wit91]。该猜想后来由 M.~Kontsevich [Kon92] 证明，它以优美的方式把理论物理、代数几何与数学物理联系在一起。近来，物理学文献中这类想法出现了复兴，与 Jackiw--Teitelboim 引力及其全息对偶 Sachdev--Ye--Kitaev 模型相关 [Kit; SSS]（参见 C.~Johnson 与 G.~Turiaci 的讲义 [Joh; Tur26]）。

% [Tutor Note: 原文把 Witten 猜想与 Kontsevich 定理的关系压缩成一句话，这里补出省略的论证：Witten 猜想断言，把 $\psi$ 类相交数组装成生成函数后，所得的 $\tau$ 函数是 KdV 层级的一个解（见定理 2.6）；Kontsevich 用矩阵模型证明了这一点——该矩阵模型的 Feynman 图恰好给出模空间上的积分。也就是说，“随机矩阵方法”不只是启发性的类比，而是猜想的证明工具，这正是本讲义 §2.3 与 §3 的主线。]

% [Tutor Note: JT（Jackiw--Teitelboim）引力的路径积分在拓扑层面归结为双曲曲面模空间上的 Weil--Petersson 体积，而后者又可由拓扑递归计算（见 §4.1 与表 3）；SYK 模型是它的全息对偶。这条线索与前半部分的 CohFT/TR 主线在结构上平行，读者可对照 §4.1 阅读。]

另一个与黎曼曲面理论有着深刻联系的物理理论是弦论。当一根弦在时空中传播时，它扫出一个黎曼曲面，即这根弦的\emph{世界面}。这些本质上就是 Feynman 图的“弦化”版本。

\figomit{1}{一根弦在时空中传播，扫出一个黎曼曲面：即弦的世界面。}

% [原文 p.3]
该理论的路径积分，在数学上被描述为从黎曼曲面到时空的映射的模空间上的积分（参见 C.-C.~M.~Liu 的讲义 [Liu]）。这类积分所满足的性质，在数学上由上同调场论这一概念来描述。

这些讲义的目标是描述与上述想法相关的数学，尤其聚焦于黎曼曲面的模空间、上同调场论的概念，以及它的递归解。

% [Tutor Note: 原文的 “recursive solution”（递归解）指的是：CohFT 的关联函数由拓扑递归（TR）给出，而 TR 只需谱曲线与初始数据，即可逐阶递推地算出全部 $\omega_{g,n}$。这是 §2.3、§3.2、§3.3 的骨架，也是本讲义除“模空间”之外的另一半内容。]

主要参考文献包括：

\begin{itemize}[leftmargin=1.5em,itemsep=0.55em,label={}]
  \item \textbf{[Zvo12]} D.~Zvonkine, ``An introduction to moduli spaces of curves and their intersection theory''
        \\[1pt] {\small\kaishu 关于曲线模空间、其相交理论以及 Witten 猜想的浅显讲义}
  \item \textbf{[Pan19]} R.~Pandharipande, ``Cohomological field theory calculations''
        \\[1pt] {\small\kaishu 关于上同调场论的浅显讲义，侧重例子}
  \item \textbf{[Sch20]} J.~Schmitt, ``The moduli space of curves''
        \\[1pt] {\small\kaishu 从代数几何取向出发、关于曲线模空间及其上同调的讲义}
  \item \textbf{[ACG11]} E.~Arbarello, M.~Cornalba, P.~A.~Griffiths, ``Geometry of Algebraic Curves, Vol.~II''
        \\[1pt] {\small\kaishu 关于黎曼曲面及其模空间的综合教材}
\end{itemize}

\medskip
\noindent{\heiti 致谢。} 我们感谢 Les Houches 暑期学校提供了非凡的工作与学习环境。我们也感谢各位学员的热情、富有洞见的问题、宝贵的意见以及令人愉快的氛围。

我们深深感谢 Ga\"etan Borot、Bertrand Eynard、Reinier Kramer、Rahul Pandharipande、Adrien Sauvaget、Sergey Shadrin 和 Dimitri Zvonkine 多年来与我们分享关于模空间的洞见。特别感谢 Sara Perletti 提供了这些讲义的基础手稿。我们还要感谢 Simon Sieroka 仔细阅读了早期草稿并帮助发现若干笔误，以及审稿人提出的建设性意见——这些意见帮助改进了最终稿件。

\medskip
\noindent{\heiti 资助信息。} 本工作部分是 ERC-SyG 项目 \emph{Recursive and Exact New Quantum Theory}（ReNewQuantum）的成果，该项目在欧盟 Horizon 2020 研究与创新计划下由欧洲研究理事会（ERC）资助，资助协议编号 810573。A.~G. 得到了 ETH Fellowship（22-2 FEL-003）以及苏黎世联邦理工学院数学研究所（Forschungsinstitut f\"ur Mathematik）的 Hermann-Weyl-Instructorship 的支持。D.~L. 得到的里雅斯特大学、INdAM 的 GNSAGA 研究组以及 INFN 的 MMNLP（APINE）项目的支持。

% [Tutor Note: 原文 “Acknowledgements” 与 “Funding information” 是正文段落（非脚注），此处照原文结构译为两段带标题的正文。“We are deeply indebted to” 译为“深深感谢”；“base manuscript”（基础手稿）指这些讲义所依据的原始笔记。]

% [Tutor Note: 原文 p.2 的 Figure 1 位于页面右下角，正文环绕排版；本译注本不复刻插图，故将图注译文置于其所在段落之后。图注原文：“A string travelling through spacetime, tracing out a Riemann surface: the worldsheet of the string.”（一根弦在时空中传播，扫出一个黎曼曲面：即弦的世界面。）]
